\chapter{La hipérbola}

\begin{center}\begin{minipage}{.88\textwidth}\small\sffamily
\textbf{Propósito.} Reconocer la hipérbola, determinar sus elementos y usar
sus asíntotas para interpretar y construir su gráfica.
\end{minipage}\end{center}

\begin{definition}{Hipérbola}{hiperbola}
Es el lugar geométrico de los puntos $P$ para los cuales el valor absoluto de
la diferencia de las distancias a dos focos es constante:
\[
|PF_1-PF_2|=2a.
\]
\end{definition}

La hipérbola tiene dos ramas, un centro, dos vértices, dos focos y dos
asíntotas. Si $a$ y $b$ son sus semiejes y $c$ es la distancia del centro a un
foco, entonces:

\begin{formula}[title={Relación fundamental}]
\[
c^2=a^2+b^2.
\]
\end{formula}

\section{Ecuaciones ordinarias}

\begin{formula}[title={Hipérbola horizontal}]
\[
\frac{(x-h)^2}{a^2}-\frac{(y-k)^2}{b^2}=1.
\]
Vértices $(h\pm a,k)$; focos $(h\pm c,k)$; asíntotas
$y-k=\pm\frac ba(x-h)$.
\end{formula}

\begin{formula}[title={Hipérbola vertical}]
\[
\frac{(y-k)^2}{a^2}-\frac{(x-h)^2}{b^2}=1.
\]
Vértices $(h,k\pm a)$; focos $(h,k\pm c)$; asíntotas
$y-k=\pm\frac ab(x-h)$.
\end{formula}

El término positivo indica la dirección en la que abren las ramas.

\begin{example}{Elementos de una hipérbola}{elementos-hiperbola}
Analiza
\[
\frac{(x-1)^2}{9}-\frac{(y+2)^2}{16}=1.
\]
\begin{solution}
\textbf{Paso 1.} $C(1,-2)$, $a=3$, $b=4$; abre horizontalmente.

\textbf{Paso 2.} $c=\sqrt{9+16}=5$.

\textbf{Paso 3.} Vértices $(-2,-2)$ y $(4,-2)$; focos $(-4,-2)$ y $(6,-2)$.

\textbf{Paso 4.} Asíntotas:
\[
y+2=\pm\frac43(x-1).
\]
\end{solution}
\end{example}

\section{Rectángulo auxiliar y asíntotas}

Para graficar se construye el rectángulo centrado en $(h,k)$, con semiancho y
semialto determinados por $a$ y $b$. Sus diagonales son las asíntotas. Las
ramas pasan por los vértices y se aproximan a ellas sin alcanzarlas.

\begin{remark}
Las asíntotas no forman parte de la hipérbola. Sirven como guías de su
comportamiento cuando los puntos se alejan del centro.
\end{remark}

\section{Excentricidad y lado recto}

\begin{formula}[title={Excentricidad de la hipérbola}]
\[
e=\frac ca>1,
\qquad
L=\frac{2b^2}{a}.
\]
\end{formula}

\begin{example}{Excentricidad}{exc-hiperbola}
Para $\frac{y^2}{25}-\frac{x^2}{144}=1$, calcula $e$ y las asíntotas.
\begin{solution}
\textbf{Paso 1.} $a=5$, $b=12$ y $c=\sqrt{25+144}=13$.

\textbf{Paso 2.} $e=\frac{13}{5}$.

\textbf{Paso 3.} Es vertical, así que
\[
y=\pm\frac abx=\pm\frac5{12}x.
\]
\end{solution}
\end{example}

\section{Obtención de la ecuación}

\begin{example}{Vértices y focos}{ecuacion-hiperbola}
Una hipérbola horizontal tiene centro $(2,1)$, vértices $(\pm3,0)+(2,1)$,
es decir, $(-1,1)$ y $(5,1)$, y focos $(-3,1)$ y $(7,1)$. Obtén su ecuación.
\begin{solution}
\textbf{Paso 1.} $a=3$ y $c=5$.

\textbf{Paso 2.} De $c^2=a^2+b^2$:
\[
b^2=25-9=16.
\]
\textbf{Paso 3.} La ecuación es
\[
\frac{(x-2)^2}{9}-\frac{(y-1)^2}{16}=1.
\]
\end{solution}
\end{example}

\section{Forma general}

En una hipérbola no rotada, los términos $x^2$ y $y^2$ tienen signos
opuestos. Se completan cuadrados y se normaliza a un lado derecho igual a $1$.

\begin{example}{De general a ordinaria}{general-hiperbola}
Analiza $4x^2-9y^2-8x-36y-68=0$.
\begin{solution}
\textbf{Paso 1.} Agrupamos:
\[
4(x^2-2x)-9(y^2+4y)=68.
\]
\textbf{Paso 2.} Completamos cuadrados:
\[
4(x-1)^2-9(y+2)^2=36.
\]
\textbf{Paso 3.} Dividimos entre $36$:
\[
\frac{(x-1)^2}{9}-\frac{(y+2)^2}{4}=1.
\]
\textbf{Paso 4.} $C(1,-2)$, $a=3$, $b=2$, $c=\sqrt{13}$; es horizontal.
\end{solution}
\end{example}

\section{Hipérbola equilátera}

Si $a=b$, las asíntotas son perpendiculares y la hipérbola se llama
equilátera o rectangular. En ejes adecuados tiene forma
\[
x^2-y^2=a^2.
\]

\section{Ejercicios y problemas}

\begin{exercise}{Hipérbolas}{ej-hiperbola}
\begin{enumerate}
  \item Analiza $\frac{x^2}{16}-\frac{y^2}{9}=1$.
  \item Analiza $\frac{(y-3)^2}{4}-\frac{(x+1)^2}{12}=1$.
  \item Obtén la ecuación con centro $(0,0)$, vértices $(\pm4,0)$ y focos
  $(\pm5,0)$.
  \item Obtén la ecuación vertical con centro $(2,-1)$, $a=3$ y $b=5$.
  \item Convierte $9x^2-4y^2+54x+8y+41=0$ a forma ordinaria.
  \item Calcula excentricidad y lado recto de $\frac{x^2}{25}-\frac{y^2}{11}=1$.
  \item Decide si $4x^2-4y^2=36$ es equilátera.
\end{enumerate}
\end{exercise}

\begin{problem}{Aplicaciones}{prob-hiperbola}
\begin{enumerate}
  \item Dos estaciones están en $(-5,0)$ y $(5,0)$. La diferencia de
  distancias desde una fuente es $6$. Obtén la ecuación del lugar geométrico.
  \item Determina las asíntotas y vértices de
  $16x^2-9y^2-64x-18y-89=0$.
  \item Encuentra las intersecciones de $y=x$ con $\frac{x^2}{9}-\frac{y^2}{16}=1$.
\end{enumerate}
\end{problem}

\section*{Resumen}\addcontentsline{toc}{section}{Resumen}
\begin{properties}
\[
\frac{(x-h)^2}{a^2}-\frac{(y-k)^2}{b^2}=1,
\qquad c^2=a^2+b^2.
\]
\[
e=\frac ca>1,
\qquad L=\frac{2b^2}{a}.
\]
El término positivo indica la orientación.
\end{properties}
